\documentclass[../Main.tex]{subfiles}
\begin{document}

\section{Conclusion}
\label{section:6.1}

This thesis developed a lightweight web-based grading application tailored for
individual teachers to support the grading of paper-based multiple-choice
examinations. By integrating a structured OMR pipeline and a two-stage OCR
pipeline based on YOLO and CRNN, the application can extract student IDs, exam
codes, selected answers, and Vietnamese handwritten names from camera-captured
answer-sheet images.

Beyond the core recognition modules, an important contribution of the thesis is
the design of a human-in-the-loop review architecture. In this workflow, the
role of OCR and OMR is not to act as a fully autonomous decision-making system,
but to assist the teacher by providing recognized outputs together with visual
evidence for verification. This allows the teacher to inspect uncertain fields,
correct recognition results when necessary, and confirm the final grading result
before it is stored.

The experimental results show that the proposed application is practically
feasible for the intended grading scenario. The OMR pipeline achieved strong
performance on the prepared structured answer-sheet dataset, while the OCR
pipeline achieved promising results for handwritten student-name recognition.
Together with the teacher-centered review workflow, these results demonstrate
that the proposed application can reduce repetitive grading work while still
preserving reliability and human control in the final grading process.

\section{Future Work}
\label{section:6.2}

Although the current application is practically useful for the intended
single-teacher grading scenario, it still leaves several directions for further
development. These directions arise naturally from the current scope of the
system, the constraints of the supported answer-sheet format, and the practical
challenges of OCR/OMR-based grading in real educational environments.

A first direction is to improve system scalability and deployment capability.
At present, the application is designed as a local single-user tool for one
teacher. Future work can extend this design into a multi-user platform with
authentication, role management, and support for multiple teachers or
school-level usage. This would make the system more suitable for broader
educational deployment and would also allow performance evaluation under
heavier and more concurrent workloads.

A second direction is to support more flexible answer-sheet layouts. Instead of
depending only on predefined region coordinates and a fixed structured layout,
future work may develop a more dynamic template-based mechanism that allows
teachers to configure or design different answer-sheet formats while still
enabling the OMR engine to adapt reliably to them. This would broaden the
practical applicability of the system beyond the current supported form type,
which is still limited to the predefined structured layout used in this thesis.

A third direction is to improve recognition and identity matching. This may
include stronger OCR models for handwritten Vietnamese names, more robust
post-processing strategies, and more advanced matching methods for uncertain
student identities. In addition to improving the current YOLO-CRNN pipeline,
future work may also explore multimodal or vision-language approaches for
particularly difficult handwritten cases. Such improvements are important
because recognition quality can still be affected by irregular handwriting,
connected strokes, diacritics, shadow, blur, and other challenging capture
conditions, which remain practical limitations of the current implementation.

A fourth direction is to improve the user-side capture workflow. For example, a
dedicated mobile application could be developed so that teachers can capture,
upload, and review answer sheets directly on smartphones in a more convenient
and integrated manner. This is meaningful because the current workflow is still
centered mainly on the web application and local browser-based interaction.

A fifth direction is to strengthen the evaluation process itself. Although the
current OCR and OMR evaluations show promising practical results, broader
testing on larger and more diverse datasets would provide stronger evidence of
generalizability. Future work can therefore include additional answer-sheet
images, more varied handwriting styles, and more difficult real-world capture
conditions in order to validate the robustness of the application more
thoroughly, since the present evaluation is still constrained by the scale and
diversity of the available project datasets.

Overall, these directions provide a practical path for extending the current
single-teacher grading application into a more flexible, robust, and scalable
educational assessment platform.

\end{document}
